\section{Summary}
\label{sec:slam-summary}

This chapter presented the design and implementation of a LiDAR-centric SLAM pipeline for GNSS-degraded operation and evaluated three configurations: the baseline GLIM system, a refined parameterised variant, and GLIM-GNSS. The results across five parameterisation datasets showed that parameter refinement reduced drift and improved trajectory smoothness relative to the baseline, while GLIM-GNSS achieved the lowest errors across all reported metrics, noting that this improvement reflects both algorithmic changes and the introduction of sparse GNSS factors as an additional sensor input unavailable to the baseline configurations. Evaluation on a held-out Evaluation Dataset confirmed that the performance hierarchy generalises to unseen routes, with GLIM-GNSS additionally demonstrating near-deterministic run-to-run consistency. The discussion highlighted characteristic failure modes in structured vineyard environments, the distinction between mean accuracy and run-to-run robustness, and the trade-off between peak and median performance arising from the parameter optimisation process.